% Class 13 Quiz and Class Activity
% Student version (no solutions).
\documentclass[11pt,a4paper]{article}

\newif\ifsolutions
\solutionsfalse

\usepackage{amsmath}
\usepackage{amssymb}
\usepackage{array}
\usepackage{caption}
\usepackage{fancyhdr}
\usepackage{float}
\usepackage{graphicx}
\usepackage{setspace}
\usepackage{needspace}
\usepackage[most]{tcolorbox}
\usepackage{hyperref}
\usepackage[margin=1in]{geometry}
\onehalfspacing
\setlength{\headheight}{14pt}
\pagestyle{fancy}
\lhead{\ifsolutions\textbf{Solutions}\else Name:\fi}
\rhead{}
\cfoot{\thepage}
\renewcommand{\labelenumii}{\textbf{\alph{enumii}.}}

\usepackage{tikz}

% ---------------------------------------------------------------------------
% Helpers
% ---------------------------------------------------------------------------
\newcommand{\classdate}{October 13, 2026}

% Title block in the same style as \maketitle (used for the activity page)
\newcommand{\classtitle}[1]{%
  \null\vskip 2em
  \begin{center}
  {\LARGE #1\par}\vskip 1.5em
  {\large PUBPOL 205\par}\vskip 1em
  {\large \classdate\par}
  \end{center}\par\vskip 1.5em}

% Answer space (student version only). The argument is a relative weight:
% \ans{2} gets twice the space of \ans{1} on the same page.
\newcommand{\ans}[1]{\ifsolutions\else\par\vspace{0pt plus #1fill}\fi}

% Page break inside a problem (student version only), with a "continued" header
\newcommand{\continuepage}[1]{%
  \ifsolutions\par\else
  \newpage\pagestyle{plain}\thispagestyle{plain}
  \noindent\textbf{#1 (continued)}\par\medskip
  \fi}

% New page before each problem (student version); the solutions flow continuously
\newcommand{\problempage}{%
  \ifsolutions\par\needspace{10\baselineskip}\bigskip
  \else\newpage\pagestyle{plain}\thispagestyle{plain}\fi}

% Solution blocks
\definecolor{solblue}{RGB}{20,70,130}
\newtcolorbox{solution}{breakable, enhanced, blanker,
  left=10pt, top=2pt, bottom=2pt, before skip=6pt, after skip=10pt,
  borderline west={2pt}{0pt}{solblue}, colupper=solblue,
  before upper={\textbf{Solution.}\ }}

% Insert a figure file if present; otherwise show a placeholder box
\newcommand{\figfile}[2]{%
  \IfFileExists{#2}{\includegraphics[width=#1]{#2}}%
  {\fbox{\parbox[c][4cm][c]{#1}{\centering\small Figure file \texttt{#2} not found}}}}

\newcommand{\MPL}{\mathit{MPL}}
\newcommand{\MPK}{\mathit{MPK}}
\newcommand{\MRTS}{\mathit{MRTS}}

% Quiz figure: isoquant, isocost line, and the cost-minimizing bundle E.
% Drawn with isoquant K = 32/L and isocost K = 8 - L/2 (W/R = 1/2), tangent at (8,4).
% Optional argument "marked" adds bundles A and B and the higher isocost through them.
\newcommand{\quizfigure}[1][]{%
\begin{tikzpicture}[x=0.45cm,y=0.55cm,font=\small]
  \draw[->] (0,0)--(19.5,0) node[below left] {Labor ($L$)};
  \draw[->] (0,0)--(0,10.8) node[above] {Capital ($K$)};
  % isocost through E
  \draw[very thick] (0,8)--(16,0) node[above right] {$C_0$};
  % isoquant
  \draw[very thick,domain=3.2:18.6,samples=120] plot (\x,{32/\x})
    node[right] {$Q=\bar{Q}$};
  % optimum
  \draw[dashed] (8,0) node[below] {$L^*$} -- (8,4) -- (0,4) node[left] {$K^*$};
  \fill (8,4) circle (3pt);
  \node[above right] at (8,4) {$E$};
  \def\tmp{#1}\def\marked{marked}%
  \ifx\tmp\marked
    \draw[dashed,thick,solblue] (0,10)--(18,1) node[right,text=solblue] {$C_1$};
    \fill[solblue] (16,2) circle (3pt) node[above right,text=solblue] {$A$};
    \fill[solblue] (4,8) circle (3pt) node[above right,text=solblue] {$B$};
  \fi
\end{tikzpicture}}

\title{Class 13 Quiz\ifsolutions\ Solutions\fi}
\author{PUBPOL 205}
\date{\classdate}

\begin{document}

%%%%%%%%%%%%%%%%%%%%%%%%%%%%%%%%%%%%%%%%%%%%%%%%%%%%%%%%%%%%%%%%%%%%%%%%%%%
% QUIZ
%%%%%%%%%%%%%%%%%%%%%%%%%%%%%%%%%%%%%%%%%%%%%%%%%%%%%%%%%%%%%%%%%%%%%%%%%%%
\maketitle
\thispagestyle{fancy}

\noindent\textbf{Notation.} $L$ = labor, $K$ = capital, $Q$ = output, $W$ = wage (price of labor), $R$ = rental rate (price of capital).
\[
\MPL = \frac{\Delta Q}{\Delta L}, \qquad \MPK = \frac{\Delta Q}{\Delta K}.
\]
The marginal product of labor ($\MPL$) is the extra output from one more unit of labor, holding capital fixed. The marginal product of capital ($\MPK$) is the extra output from one more unit of capital, holding labor fixed.

\begin{enumerate}
\item In words, define the marginal rate of technical substitution of labor for capital, $\MRTS_{LK}$. What does it measure? What would $\MRTS_{LK} = 2$ mean for the firm?
\ans{1}

\end{enumerate}

\noindent\textbf{Use the following.} The marginal rate of technical substitution equals the ratio of marginal products, $\MRTS_{LK} = \MPL/\MPK$. A firm minimizes the cost of producing a given output at the bundle where the isoquant is tangent to the lowest isocost line (bundle $E$ in the figure):
\[
\underbrace{-\frac{W}{R}}_{\text{slope of isocost}} = \underbrace{-\frac{\MPL}{\MPK}}_{\text{slope of isoquant}}
\quad\Longrightarrow\quad
\frac{\MPL}{W} = \frac{\MPK}{R}.
\]
\begin{center}
\quizfigure\\[2pt]
{\small Isoquant $Q=\bar{Q}$ and isocost line $C_0$ (slope $-W/R$). $E$ is the cost-minimizing bundle.}
\end{center}

\continuepage{Class 13 Quiz}

\begin{enumerate}\setcounter{enumi}{1}
\item Suppose that at the firm's current bundle, $\dfrac{\MPK}{R} > \dfrac{\MPL}{W}$.

\noindent\textbf{a.} What does this inequality mean? Is the firm minimizing cost? Should it use more capital or more labor? Explain.
\ans{3}


\noindent\textbf{b.} In this case, is the slope of the isocost line (the firm's budget constraint) bigger or smaller than $\MRTS_{LK}$? Compare absolute values: $W/R$ versus $\MRTS_{LK}$. Then, on the figure on page 1, spot a bundle on the isoquant where this is true and label it $A$.
\ans{2}


\item Suppose instead that $\dfrac{\MPL}{W} > \dfrac{\MPK}{R}$.

\noindent\textbf{a.} What does this inequality mean? Is the firm minimizing cost? Should it use more capital or more labor? Explain.
\ans{3}


\noindent\textbf{b.} In this case, is the slope of the isocost line bigger or smaller than $\MRTS_{LK}$? Compare absolute values: $W/R$ versus $\MRTS_{LK}$. Then, on the figure on page 1, spot a bundle on the isoquant where this is true and label it $B$.
\ans{2}

\end{enumerate}

%%%%%%%%%%%%%%%%%%%%%%%%%%%%%%%%%%%%%%%%%%%%%%%%%%%%%%%%%%%%%%%%%%%%%%%%%%%
% CLASS ACTIVITY
%%%%%%%%%%%%%%%%%%%%%%%%%%%%%%%%%%%%%%%%%%%%%%%%%%%%%%%%%%%%%%%%%%%%%%%%%%%
\newpage
\pagestyle{fancy}
\classtitle{Class 13 Activity\ifsolutions\ Solutions\fi}
\thispagestyle{fancy}

\begin{enumerate}
\item \textbf{Funnel cake costs.}

You are the CEO of a major producer of funnel cakes. Your cost accountant has provided you with a table describing your cost structure, but you have inadvertently dripped cooking grease on it and most of the table is illegible. Reconstruct the table below, given the remaining legible numbers. Show your calculations below the table.
\begin{center}
\begingroup\singlespacing\renewcommand{\arraystretch}{1.9}
\begin{tabular}{|c|*{7}{>{\centering\arraybackslash}p{1.3cm}|}}
\hline
$Q$ & $TC$ & $FC$ & $VC$ & $MC$ & $AVC$ & $AFC$ & $ATC$ \\ \hline
0 &     &  &    & -- & -- & -- & -- \\ \hline
1 &     &  & 17 &    &    &    &    \\ \hline
2 &     &  &    & 15 &    &    &    \\ \hline
3 & 101 &  &    &    &    &    &    \\ \hline
4 &     &  &    &    & 14.5 &  &    \\ \hline
5 & 122 &  & 67 &    &    & 11 &    \\ \hline
6 &     &  &    &    &    &    & 21 \\ \hline
\end{tabular}\endgroup
\end{center}
\ans{1}


\problempage
\item \textbf{Cost minimization and the cost function.}

Suppose a firm has the following production function: $Q = 2KL$. The marginal product of capital for this production function is $2L$, and the marginal product of labor is $2K$. If capital rents for \$100 per unit per day, labor can be hired for \$200 per unit per day, and the firm is minimizing costs,

\noindent\textbf{a.} Find the total cost of producing 60 units of output.
\ans{2.2}


\noindent\textbf{b.} Suppose you want to find the cost of producing a more general $q$ units of output. Using the same set of procedures you followed in (a), but using a more general $q$ instead of 60, find this firm's total cost function, which expresses total cost as a function of $q$.
\ans{2.2}


\noindent\textbf{c.} Divide the total cost function you found in (b) by $q$ to find this firm's average cost function.
\ans{1}


\problempage
\item \textbf{Hack's Berries.}

Hack's Berries faces a short-run total cost of production given by
\[
TC = Q^3 - 12Q^2 + 100Q + 1{,}000,
\]
where $Q$ is the number of crates of berries produced per day. Hack's marginal cost of producing berries is $MC = 3Q^2 - 24Q + 100$.

\noindent\textbf{a.} What is the level of Hack's fixed cost?
\ans{1}


\noindent\textbf{b.} What is Hack's short-run average variable cost of producing berries?
\ans{1.3}


\noindent\textbf{c.} If berries sell for \$60 per crate, how many berries should Hack produce? How do you know? (\emph{Hint:} You may want to remember the relationship between $MC$ and $AVC$ when $AVC$ is at its minimum.)
\ans{3}


\continuepage{3. Hack's Berries}

\noindent\textbf{d.} If the price of berries is \$73 per crate, how many berries should Hack produce? Explain.
\ans{1}


\problempage
\ifsolutions\needspace{0.45\textheight}\fi
\item \textbf{Potato farming and interest rates.}

Consider the graph below, which depicts the cost curves of a perfectly competitive seller of potatoes. Potatoes currently sell for \$3 per pound.
\begin{center}
\figfile{0.55\textwidth}{c13_potato_costs.jpeg}
\end{center}

\noindent\textbf{a.} To maximize profit, how many pounds of potatoes should this seller produce?
\ans{1}


Suppose that the potato grower's bank ratchets up the interest rate applicable to the grower's adjustable-rate mortgage loan. This increases the size of the potato grower's monthly mortgage payment.

\noindent\textbf{b.} Illustrate the change in the mortgage payment by shifting the appropriate cost curves on the graph above.
\ans{0.6}


\continuepage{4. Potato farming and interest rates}

\noindent\textbf{c.} Which curves shift? Which do not? Why?
\ans{1}


\noindent\textbf{d.} How does the change in interest rates affect the grower's decision on how many potatoes to produce?
\ans{1}


\noindent\textbf{e.} What happens to the potato grower's profit as a result of the increased interest rate?
\ans{0.8}


\noindent\textbf{f.} How does the change in interest rates affect the shape and/or position of the grower's short-run supply curve?
\ans{1}

\end{enumerate}

\end{document}
